% Volume 1, Chapter 2: Implementation and Specification Guidelines
\chapter{IEEE 752 Implementation and Specification Guidelines}
\label{ch:ieee752-spec-guidelines}

\section{Specification Guidelines}

\subsection{Schema Compliance}

\begin{itemize}
    \item \textbf{Protobuf}: Use \texttt{uint64}/\texttt{int64} for timestamp, count, index, and size fields. No \texttt{uint32}/\texttt{int32} for such geometrica fields.
    \item \textbf{JSON Schema}: Use \texttt{"type": "integer"} with \texttt{"minimum"}/\texttt{"maximum"} or format hint for 64-bit range where applicable.
    \item \textbf{Avro}: Use \texttt{long} for 64-bit integers; \texttt{double} for 64-bit float.
\end{itemize}

\subsection{Implementation Checklist}

When implementing data structures that participate in the Z-field:

\begin{enumerate}
    \item Identify all numeric fields that represent timestamps, counts, indices, or sizes.
    \item Map each to \texttt{UInt64}, \texttt{Int64}, or \texttt{Double} (IEEE 754) per the language mapping.
    \item Ensure wire formats (e.g.\ Gated Loopback Flower \texttt{created\_at\_epoch\_ms}) use \texttt{uint64}.
    \item Document any range limits where the runtime does not provide native 64-bit types.
\end{enumerate}

\section{References}

\begin{itemize}
    \item Gated Loopback Flower: \texttt{created\_at\_epoch\_ms} is \texttt{uint64} (compliant).
    \item Node identity hex space: \(0x00\)--\(0x63\) (0x64 slots) is separate from the 64-bit Z-field; both apply.
    \item IEEE 754: ISO/IEC/IEEE 60559 (binary64) for double-precision floating-point when used.
\end{itemize}
